\chapter{实验结果与分析}

\section{引言}

第 3 章详细介绍了 STCA 模型的算法设计。本章将在外域场景下评估模型的性能，验证模型的跨场景泛化能力。首先介绍实验设置与评价指标，然后展示实验结果并进行分析，最后通过消融实验验证各核心模块的有效性。

\section{实验设置}

\subsection{硬件环境}
实验在 Intel Core i7 处理器、NVIDIA GeForce RTX 3060 GPU、16 GB 内存的工作站上完成。模型基于 PyTorch 框架，使用 CUDA 11.8 和 cuDNN 加速计算。

\subsection{超参数配置}
实验超参数配置见表~\ref{tab:hyperparameters}。训练使用 Adam 优化器，学习率在训练集 loss 停滞时自动衰减。损失函数使用二元交叉熵（BCEWithLogitsLoss），内部集成 Sigmoid 激活，数值稳定性更好。

\begin{table}[h]
\centering
\caption{模型超参数配置}
\label{tab:hyperparameters}
\begin{tabular}{lc}
\hline
参数 & 取值 \\
\hline
随机种子 & 42 \\
训练轮数 & 50 \\
批次大小 & 16 \\
初始学习率 & $10^{-4}$ \\
空间嵌入维度 & 16 \\
时序嵌入维度 & 16 \\
交叉注意力嵌入维度 & 32 \\
稀疏嵌入维度 & 32 \\
Dropout 率 & 0.5 \\
\hline
\end{tabular}
\end{table}

\section{实验结果}

外域泛化实验采用跨测点划分策略：5 个测点（P2、P3、P4、P5、P8）的数据用于训练，2 个测点（P6、P7）的数据用于测试。该策略模拟真实部署场景：模型在部分地点训练后，需泛化到未见过的地点。

\subsection{训练过程分析}

图~\ref{fig:training_history_outdomain} 展示了外域场景下的训练历史。训练集损失随训练轮次增加平滑下降，在约 35 轮后趋于稳定，最终训练集 Loss 收敛至 0.5821。训练过程中未出现明显过拟合现象，表明 Dropout 正则化策略有效。

\begin{figure}[htbp]
\centering
\includegraphics[width=0.85\textwidth]{figures/training_history_outdomain.png}
\caption{外域场景下训练历史曲线}
\label{fig:training_history_outdomain}
\end{figure}

\subsection{混淆矩阵分析}

图~\ref{fig:confusion_matrix_outdomain} 为外域测试集上的混淆矩阵结果。模型对 LOS 信号的识别正确率为 79.52\%，对 NLOS 信号的识别正确率为 72.32\%。总体准确率为 76.02\%，表明模型在跨场景下具有一定的泛化能力。

\begin{figure}[htbp]
\centering
\includegraphics[width=0.6\textwidth]{figures/confusion_matrix_outdomain.png}
\caption{外域测试集混淆矩阵}
\label{fig:confusion_matrix_outdomain}
\end{figure}

\subsection{ROC 曲线与 PR 曲线}

图~\ref{fig:curves_outdomain} 为外域测试集上的 ROC 曲线和 PR 曲线。ROC-AUC 为 84.07\%，PR-AUC 为 83.71\%，表明模型在不同决策阈值下均保持稳定的类别分离能力。ROC 曲线快速上升至左上角区域，PR 曲线在高精确率区间维持较长时间，反映了模型对 NLOS 信号的识别稳定性。

\begin{figure}[htbp]
\centering
\includegraphics[width=0.85\textwidth]{figures/roc_curve_outdomain.png}
\caption{外域测试集 ROC 曲线}
\label{fig:curves_outdomain}
\end{figure}

\begin{figure}[htbp]
\centering
\includegraphics[width=0.85\textwidth]{figures/pr_curve_outdomain.png}
\caption{外域测试集 PR 曲线}
\label{fig:pr_curve_outdomain}
\end{figure}

\subsection{定量评估指标}

表~\ref{tab:metrics_outdomain} 汇总了模型在外域测试集上的各项性能指标。

\begin{table}[h]
\centering
\caption{外域测试集分类性能指标}
\label{tab:metrics_outdomain}
\begin{tabular}{cccccc}
\hline
准确率 & 精确率 & 召回率 & F1 分数 & ROC-AUC & PR-AUC \\
\hline
76.02\% & 78.79\% & 72.32\% & 75.42\% & 84.07\% & 83.71\% \\
\hline
\end{tabular}
\end{table}

从应用角度看，高召回率对 NLOS 检测更为关键：漏检 NLOS 信号可能在定位解算中引入较大误差，而误检 LOS 信号仅减少可用观测数量。本模型在保持合理精确率的同时实现了较好的召回率，满足实际应用需求。

\section{消融实验}

为验证 STCA 模型各组成部分的有效性，本研究设计了系统的消融实验。消融实验分为两部分：时间窗口大小敏感性分析和模块消融实验。

\subsection{时间窗口大小敏感性分析}

时间窗口大小决定了时序编码器接收的历史信息长度，直接影响模型对信号时序演化模式的捕捉能力。窗口过小则历史信息不足，难以捕捉 NLOS 状态的时序连续性；窗口过大则引入冗余噪声，增加计算负担且易导致过拟合。

本研究测试了 14 种不同的窗口大小配置，分别为 6、8、10、12、14、16、18、20、22、24、26、28、30、32 历元。表~\ref{tab:ablation_window} 展示了不同窗口大小下的模型性能对比结果。

\begin{table}[h]
\centering
\caption{时间窗口大小消融实验结果}
\label{tab:ablation_window}
\begin{tabular}{cccc}
\hline
窗口大小 & 准确率 & 精确率 & 召回率 \\
\hline
6 & 79.34\% & 72.02\% & 83.35\% \\
8 & 82.12\% & 74.93\% & 86.56\% \\
10 & 80.61\% & 72.63\% & 86.76\% \\
12 & 80.44\% & 73.76\% & 83.33\% \\
14 & 81.45\% & 75.62\% & 82.76\% \\
16 & 79.53\% & 73.59\% & 80.41\% \\
18 & 81.29\% & 75.50\% & 82.52\% \\
20 & 80.07\% & 73.62\% & 82.42\% \\
22 & 80.31\% & 72.35\% & 86.55\% \\
24 & 78.88\% & 71.43\% & 83.52\% \\
26 & 79.43\% & 72.22\% & 83.60\% \\
28 & 80.29\% & 72.46\% & 86.25\% \\
30 & 78.94\% & 71.27\% & 84.30\% \\
32 & 79.76\% & 71.03\% & 88.24\% \\
\hline
\end{tabular}
\end{table}

实验结果表明，窗口大小为 8 时模型性能最优，准确率达到 82.12\%，F1 分数达到 80.33\%。窗口小于 8 时，随着窗口增大，模型能够利用更多的历史时序信息，性能逐步提升；窗口大于 8 时，继续增加窗口长度反而导致性能下降，原因是过多的历史信息引入了冗余噪声，且增加了模型的学习难度。因此，后续实验均采用窗口大小 8 作为默认配置。

\begin{figure}[htbp]
\centering
\includegraphics[width=0.9\textwidth]{figures/ablation_study_results.png}
\caption{消融实验结果对比图}
\label{fig:ablation_results}
\end{figure}

\subsection{模块消融实验}

为定量分析 STCA 模型各功能模块对整体性能的贡献，本研究设计了模块消融实验。实验采用控制变量法，在保持其他条件不变的情况下，依次移除或替换特定模块，观察模型性能变化。消融实验算法如算法~\ref{alg:ablation} 所示。

\begin{algorithm}[htbp]
\caption{STCA 模型消融实验算法}
\label{alg:ablation}
\begin{algorithmic}[1]
\Require 测试集 $\mathcal{D}_{test}$, 模块配置列表 $\mathcal{C} = \{c_1, c_2, \ldots, c_k\}$
\Ensure 各配置的Performance 指标 $\mathcal{M} = \{m_1, m_2, \ldots, m_k\}$

\State $\mathcal{M} \gets \emptyset$
\For{每个配置 $c_i \in \mathcal{C}$}
    \State 加载对应配置的模型参数 $\Theta_i$
    \State $\text{TP}, \text{TN}, \text{FP}, \text{FN} \gets 0$
    \For{$(\bm{X}_{spatial}^{(j)}, \bm{X}_{temporal}^{(j)}, y_j) \in \mathcal{D}_{test}$}
        \State $\hat{y}_j \gets \text{STCA\_Forward}(\bm{X}_{spatial}^{(j)}, \bm{X}_{temporal}^{(j)}; \Theta_i)$
        \State $\hat{label}_j \gets \mathbb{I}(\hat{y}_j > 0.5)$ \Comment{阈值判定}
        \If{$\hat{label}_j = y_j$}
            \If{$y_j = 1$} $\text{TP} \gets \text{TP} + 1$
            \Else $\text{TN} \gets \text{TN} + 1$
            \EndIf
        \Else
            \If{$y_j = 1$} $\text{FN} \gets \text{FN} + 1$
            \Else $\text{FP} \gets \text{FP} + 1$
            \EndIf
        \EndIf
    \EndFor
    \State $\text{Acc}_i \gets \frac{\text{TP} + \text{TN}}{\text{TP} + \text{TN} + \text{FP} + \text{FN}}$
    \State $\text{Precision}_i \gets \frac{\text{TP}}{\text{TP} + \text{FP}}$
    \State $\text{Recall}_i \gets \frac{\text{TP}}{\text{TP} + \text{FN}}$
    \State $\text{F1}_i \gets \frac{2 \times \text{Precision}_i \times \text{Recall}_i}{\text{Precision}_i + \text{Recall}_i}$
    \State $\mathcal{M} \gets \mathcal{M} \cup \{(\text{Acc}_i, \text{F1}_i)\}$
\EndFor
\State \Return $\mathcal{M}$
\end{algorithmic}
\end{algorithm}

实验分为三组：
（1）基线模型（Baseline）：移除交叉注意力模块和稀疏表示层，仅保留空间编码器和时序编码器，融合方式采用简单的特征拼接；
（2）+交叉注意力（+CrossAttn）：在基线模型基础上添加交叉注意力融合模块，时序特征为 Query，空间特征为 Key/Value；
（3）完整模型（Full STCA）：包含空间编码器、时序编码器、交叉注意力模块和稀疏表示层的完整架构。

表~\ref{tab:ablation_modules} 展示了各配置的性能对比结果。

\begin{table}[h]
\centering
\caption{模块消融实验结果}
\label{tab:ablation_modules}
\begin{tabular}{lcc}
\hline
模型配置 & 准确率 & F1 分数 \\
\hline
基线模型（无交叉注意力 + 稀疏表示） & 75.23\% & 73.85\% \\
+交叉注意力 & 79.87\% & 77.92\% \\
完整 STCA 模型 & 82.12\% & 80.33\% \\
\hline
\end{tabular}
\end{table}

实验结果表明，每个模块都对最终性能有正向贡献：
（1）基线模型准确率为 75.23\%，验证了简单的特征拼接无法充分建模时空特征的复杂交互关系；
（2）添加交叉注意力模块后，准确率提升至 79.87\%，表明交叉注意力机制能够有效建立时空特征之间的自适应关联；
（3）完整 STCA 模型准确率达到 82.12\%，较基线模型提升 6.89 个百分点，验证了交叉注意力模块和稀疏表示层的有效性。

综上所述，STCA 模型的时空双通道架构能够同时利用空间分集和时间相关性信息，交叉注意力机制实现了时空特征的自适应融合，稀疏表示层增强了特征的判别能力。

\section{本章小结}

本章在外域场景下评估了 STCA 模型的性能。实验结果表明，模型在外域测试集上准确率达到 76.02\%，ROC-AUC 为 84.07\%，表明模型在跨场景下具有一定的泛化能力。时间窗口消融实验表明窗口大小为 8 时性能最优（准确率 82.12\%，F1 分数 80.33\%）。模块消融实验验证了交叉注意力模块和稀疏表示层对整体性能的贡献，完整模型较基线模型提升 6.89 个百分点。
